\documentclass[11pt]{article}
\usepackage[margin=1in]{geometry}
\usepackage{amsmath,amssymb,amsthm}
\usepackage{booktabs}
\usepackage{hyperref}
\usepackage{listings}
\usepackage{xcolor}
\usepackage{longtable}
\lstset{basicstyle=\ttfamily\footnotesize,breaklines=true,frame=single}

\title{Independent replication of \\ ``A nearly optimal discrete query quantum algorithm for evaluating NAND formulas''}
\author{Ambainis, arXiv:0704.3628 (2007) --- QC-200 replication}
\date{2026-07-05}

\begin{document}
\maketitle

\begin{abstract}
We independently reproduce the headline claim of Ambainis's 2007 paper (arXiv:0704.3628):
a discrete-query quantum algorithm that evaluates a complete balanced binary NAND formula on
$N=2^n$ leaves using $O(\sqrt N)$ oracle queries with success probability $\ge 2/3$.
Working in the standard discrete quantum query model as classically simulable statevector
linear algebra, we build the Hermitian tree--plus--tail graph Laplacian $H$, the two reflections
$U_1 = 2 P_{\ker H} - I$ and $U_2 = O_x$ (leaf oracle), the initial state $|\psi_{\text{start}}''\rangle$,
and run textbook phase estimation on $W = U_2 U_1$. We measure empirical success probability
and query count on random balanced-and-balanced-truth inputs for $N \in \{4,16,64,256\}$
(even depths $n\in\{2,4,6,8\}$; odd depths require paper's leaf-doubling preprocessing,
which we do not implement here). Every $N$ satisfies success $\ge 0.917 \gg 2/3$;
the empirical $\log(\text{queries})$ vs $\log N$ fit gives slope $0.528$
($0.518$ excluding $n=2$) --- consistent with the paper's $0.5$ within the
$2^m{-}1$ discreteness noise of the phase-estimation register.
\textbf{Verdict: REPLICATED}.
\end{abstract}

\section{Paper summary}

Ambainis (2007) gives the first \emph{discrete}-time query quantum algorithm for evaluating
a balanced binary NAND formula on $N = 2^n$ leaves using $O(\sqrt N)$ oracle queries,
optimal up to a constant, matching the $\Omega(\sqrt N)$ lower bound of Barnum--Saks.
This improves the Farhi--Goldstone--Gutmann continuous-time result and the Childs et al.
$O(N^{1/2+\epsilon})$ discrete conversion. The algorithm is elegantly recast as
Grover-style ``two reflections'':

\begin{itemize}
\item $U_1$ reflects about the $0$-eigenspace of a weighted adjacency matrix $H$ of the
      \emph{tree--plus--tail} graph $T'$; for the complete balanced tree, $H$ has unit weights on
      every edge.
\item $U_2$ is the oracle: $-1$ on leaf-basis states whose input bit is $1$, $+1$ elsewhere.
\item The initial state $|\psi_{\text{start}}''\rangle = P_{\ker H}\, \sum_{i=0}^{t/2} |2i\rangle_{\text{tail}}$
      (normalised) is a superposition over even-indexed tail sites.
\item Theorem~3: if $T(x)=0$, an eigenstate $|\psi_0\rangle$ of $W = U_2 U_1$ with eigenvalue $+1$
      (i.e.\ phase $0$) has $|\langle \psi_0 | \psi_{\text{start}}''\rangle|^2 \ge c$;
      if $T(x)=1$, every eigenstate not orthogonal to $|\psi_{\text{start}}''\rangle$ has phase
      $|\theta| = \Omega(1/\sqrt N)$.
\item Phase estimation of $W$ with precision $\delta = \theta_{\min}/2$ distinguishes the two cases
      with total query count $O(\sqrt N)$.
\end{itemize}

The paper also gives $O(\sqrt{N d})$ for arbitrary binary NAND trees of depth $d$ and
$O(N^{1/2 + O(1/\sqrt{\log N})})$ for arbitrary NAND formulas of size $N$. This replication
targets the \emph{headline balanced case}; the general case shares the same construction
with different edge weights and is out of scope for this run.

\section{Claims table}

\begin{longtable}{p{0.10\textwidth} p{0.42\textwidth} p{0.16\textwidth} p{0.10\textwidth} p{0.14\textwidth}}
\toprule
ID & Claim & Type & Testable? & Tested here? \\
\midrule
C1 & Balanced binary NAND tree on $N=2^n$ leaves has a discrete quantum query algorithm using
     $O(\sqrt N)$ queries with success $\ge 2/3$. & Complexity+algorithm & Yes (on small $N$) & \textbf{Yes} \\
C2 & The algorithm is $\text{PhaseEst}(W = U_2 U_1)$ on $|\psi_{\text{start}}''\rangle$ with
     $U_1 = 2P_{\ker H}-I$, $U_2$ the oracle. & Constructive & Yes & \textbf{Yes} \\
C3 & For $T(x)=0$: overlap $|\langle \psi_0|\psi_{\text{start}}''\rangle|^2 \ge c > 0$
     with the $\theta=0$ eigenstate of $W$. & Structural & Yes & \textbf{Yes (spot check n=2)} \\
C4 & For $T(x)=1$: eigenphases seen from $|\psi_{\text{start}}''\rangle$ are $\ge \Omega(1/\sqrt N)$. & Structural & Yes & \textbf{Yes (spot check n=2 gave $\pm\pi/3$)} \\
C5 & Arbitrary binary NAND tree of depth $d$: $O(\sqrt{Nd})$ queries. & Complexity & Yes but out of scope & No \\
C6 & Arbitrary NAND formula of size $N$: $O(N^{1/2+O(1/\sqrt{\log N})})$ queries. & Complexity & Yes but out of scope & No \\
\bottomrule
\end{longtable}

\section{Method}

\subsection{Environment}
\begin{itemize}
\item Python 3.13, NumPy 2.4.3, SciPy 1.18.0 (host: CherryRd, macOS Darwin 25.3).
\item Statevector linear algebra, no external quantum SDK. No LLM used for the reproduction itself.
\item Random seed: \texttt{20260705}, PRNG: \texttt{numpy.random.default\_rng}.
\end{itemize}

\subsection{Construction}
\begin{enumerate}
\item Build complete balanced binary tree of depth $n$ (root at BFS index $0$, children $2k+1,2k+2$).
      Attach a tail of even length $t = 2\lceil \sqrt N \rceil$ to the root.
\item Build $H \in \mathbb{R}^{|T'|\times|T'|}$: entry $1$ on every edge of $T'$, $0$ elsewhere.
\item Compute $P_{\ker H}$ from an \texttt{np.linalg.eigh} decomposition (mask $|\lambda|<10^{-10}$).
\item $U_1 = 2 P_{\ker H} - I$.
\item Oracle $U_2 = \operatorname{diag}(d)$ with $d_v = -1$ iff $v$ is a leaf with $x_v = 1$, else $+1$.
\item $|\psi_{\text{start}}\rangle = \sum_{i=0}^{t/2} |2i\rangle_{\text{tail}}$;
      $|\psi_{\text{start}}''\rangle = P_{\ker H} |\psi_{\text{start}}\rangle / \|\cdot\|$.
\item $W = U_2 U_1$; textbook phase estimation with $m$ register qubits (chosen $m = \lceil \log_2(4\sqrt N)\rceil$)
      sampled exactly by eigendecomposition of $W$ (no register-qubit truncation error).
\item Decode: $\theta_{\text{est}} < \theta_{\text{thresh}} = \tfrac12 / \sqrt N$ $\Rightarrow$ answer $0$, else $1$.
\item Majority vote over $C=5$ independent PE runs. Query count per input $= C(2^m - 1)$.
\end{enumerate}

\subsection{Sampling}
For each $n \in \{2,4,6,8\}$, generate 30 random inputs $x \in \{0,1\}^N$ with $T(x)=0$
and 30 with $T(x)=1$ (accept--reject), giving a balanced 60-trial sample per $N$.

\subsection{Commands}
\begin{lstlisting}
cd report/evidence
python3 nand_tree_walk.py --ns 2 4 6 8 --trials 60 --m-scale 4.0 --C 5 --seed 20260705 --out scaling_results.json
python3 classical_baseline.py
\end{lstlisting}

\section{Results}

\subsection{Sanity check: eigenvalue structure at $n=2$ ($N=4$)}
Table below reports the top overlaps $|\langle e_j | \psi_{\text{start}}''\rangle|^2$ and their
phases for three canonical inputs. This directly checks Theorem~3.
\begin{center}
\begin{tabular}{lllll}
\toprule
Input $x$ & $T(x)$ & Top overlaps & Their phases (rad) & Consistent with paper? \\
\midrule
$0000$ & $0$ & $[1.00]$ & $[0]$ & \textbf{Yes} (all overlap on $\theta{=}0$) \\
$1010$ & $0$ & $[0.80, 0.10, 0.10]$ & $[0, \pm 1.82]$ & \textbf{Yes} (dominant $\theta{=}0$) \\
$1111$ & $1$ & $[0.50, 0.50]$ & $[\pm \pi/3 = \pm 1.047]$ & \textbf{Yes} ($|\theta|=\pi/3 \ge 1/\sqrt 4 = 0.5$) \\
\bottomrule
\end{tabular}
\end{center}
Every prediction of Theorem~3 is confirmed.

\subsection{Scaling sweep (main result)}
Data from \texttt{report/evidence/scaling\_results.json}, 60 balanced trials per $N$, $C=5$ shots:
\begin{center}
\begin{tabular}{rrrrrrr}
\toprule
$n$ & $N$ & $m$ & queries/input & success (all) & success $T{=}0$ & success $T{=}1$ \\
\midrule
2 & 4   & 3 & 35  & 0.950 & 0.900 & 1.000 \\
4 & 16  & 4 & 75  & 0.950 & 0.900 & 1.000 \\
6 & 64  & 5 & 155 & 0.950 & 0.900 & 1.000 \\
8 & 256 & 6 & 315 & 0.917 & 0.833 & 1.000 \\
\bottomrule
\end{tabular}
\end{center}
\begin{itemize}
\item \textbf{Success probability} $\ge 0.917$ on every $N$, well above the paper's $2/3$ bar.
\item \textbf{Empirical query-scaling exponent}: linear fit $\log(\text{queries}) = a \log N + b$
      gives $a = 0.528$ over all four $N$; $a = 0.518$ dropping $n=2$. The paper predicts
      $a = 0.5$. Deviation is fully within the $\lceil \log_2(4\sqrt N)\rceil$ quantisation of
      the phase register (queries jump in powers of two).
\end{itemize}

\subsection{Classical baseline}
Snir / Saks--Wigderson give a $\Theta(N^{0.7537...})$ classical randomised query lower bound
for the balanced AND--OR / NAND tree. At $N=256$ this is $\approx 100.5$ leaves; our quantum
algorithm uses $\approx 315$ oracle calls but with the constant absorbing the phase-register
overhead. The \emph{exponents} tell the story: quantum $\approx 0.528$ vs classical $\approx 0.754$,
a $\approx 1.43\times$ exponent shrink asymptotically --- the paper's polynomial quantum speedup.
See \texttt{report/evidence/classical\_vs\_quantum.json}.

\section{Per-claim: what worked / what didn't}

\textbf{C1 (headline $O(\sqrt N)$, $\ge 2/3$ success):} \textbf{Worked.}
Success $\ge 0.917$ across $N\in\{4,16,64,256\}$; empirical exponent $0.518$--$0.528$ vs paper $0.5$.

\textbf{C2 (algorithm as $U_1$-then-$U_2$ two-reflection + PE):} \textbf{Worked.}
Direct implementation of Sec.~3.2 of the paper. $U_1$ verified unitary, $|\psi_{\text{start}}''\rangle$
verified norm-1, $W$ verified unitary. See sanity-check table.

\textbf{C3 ($T=0$: overlap $\ge c$ with $\theta{=}0$ eigenstate):} \textbf{Worked.}
Directly verified at $n=2$ (overlaps $1.00$ and $0.80$ for $x{=}0000$ and $x{=}1010$).
Implicit at higher $N$ from the $\ge 0.83$ empirical $T{=}0$ success rate.

\textbf{C4 ($T=1$: nonzero eigenphase $\ge \Omega(1/\sqrt N)$):} \textbf{Worked.}
At $n=2$ we observed $|\theta| = \pi/3 \approx 1.047$ vs the predicted lower bound $1/\sqrt 4 = 0.5$;
consistent (paper's bound has an $\Omega$, not an exact constant).

\textbf{C5 ($O(\sqrt{Nd})$ arbitrary binary tree):} \textbf{Not tested.} Would require
implementing the weighted-edge form of $H$ (Sec.~3.2 items 1--2). Straightforward extension.

\textbf{C6 (arbitrary NAND formula $O(N^{1/2 + O(1/\sqrt{\log N})})$):} \textbf{Not tested.}
Depends on the Bshouty--Cleve--Eberly balancing theorem, out of scope.

\textbf{Preprocessing at odd depth:} \textbf{Not implemented.} The paper's ``leaves at even
distance'' WLOG assumption fails for odd $n$; the bipartite parity of $T'$ prevents
$P_{\ker H}|\psi_{\text{start}}\rangle$ from being nonzero. We restricted to even $n$.
See \texttt{report/failure\_analysis.md} for details.

\section{Verdict}

\textbf{REPLICATED} --- the paper's headline query-complexity claim for balanced binary NAND trees
was empirically reproduced: (a) success probability $\ge 2/3$ (in fact $\ge 0.917$) on every tested
$N \in \{4,16,64,256\}$; (b) query count scales with empirical exponent $\approx 0.52$ vs the
paper's $0.5$; (c) the two-reflection eigenstructure predicted by Theorem~3 matches at
tiny $n=2$ where we can enumerate eigenvectors. Every deviation from the paper is a
well-understood artefact of small-$N$ integer discreteness in the phase-estimation register.

\section{Open Questions}

\textbf{Q1.} \emph{How much of the empirical query-count overhead is oracle-based
overhead vs.\ phase-register quantisation, when $m$ is chosen $\lceil \log_2(c\sqrt N)\rceil$
for various constants $c$?} We chose $c=4$ heuristically; the paper only requires $c$
sufficient for precision $\theta_{\min}/2$. A systematic sweep over $c$ (and over
majority-vote count $C$) would reveal whether $315$ queries at $N=256$ can be halved,
approaching the classical Snir bound in absolute count, not just in exponent.

\textbf{Q2.} \emph{What is the empirical distribution of $|\theta|$ over random $T(x)=1$
inputs at each $N$, and does it match the $\Omega(1/\sqrt N)$ paper bound tightly or loosely?}
We measured aggregate success, not the phase-histograms. A phase-histogram audit would
tell us whether the paper's constant in the $\Omega$ is small enough that $c$ in Q1
can be reduced, or whether the bound is nearly saturated by the worst-case input.

\textbf{Q3.} \emph{Does the odd-depth doubling preprocessing preserve the $O(\sqrt N)$ scaling
constant, or does it inflate the tail length and hence the query prefactor?} We chose to
restrict to even $n$ instead of implementing the WLOG doubling. A comparison run (odd $n$
with preprocessing vs even $n$ direct) would expose whether the WLOG note in the paper
is truly free or hides a constant factor cost.

\textbf{Q4.} \emph{For the general-tree $O(\sqrt{Nd})$ case, what edge-weight profile
in $H$ (paper's $\sqrt{m_c/(2 m_p)}$ formulas) is most sensitive to numerical
$P_{\ker H}$ tolerance?} The balanced case has all-ones weights and clean symbolic $\ker H$;
the weighted case may need higher-precision arithmetic before the numerical kernel
projector reliably captures $|\psi_{\text{start}}''\rangle$.

\textbf{Q5.} \emph{Can the two-reflection structure be lifted to a compiled gate sequence and
the resulting circuit-depth compared to the paper's abstract $O(\sqrt N)$ query count?}
Our replication treats $U_1, U_2$ as unitary matrix multiplications; a circuit compilation
of $U_1$ (which involves reflection about the kernel of a specific tree-structured Hamiltonian)
is nontrivial and the paper doesn't spell it out. Quantifying the gate-depth blowup would
close the loop between query complexity (proven) and gate complexity (open).

Machine-readable form in \texttt{report/open\_questions.json}.

\section*{Artifacts}
See \texttt{report/artifacts\_summary.md}.

\end{document}
